% ============================================================
%  ch03b：§3 谱 Eisenstein 级数：3.3 谱 Eisenstein 级数、
%         3.4 谱 vs. 函数论 Eisenstein 级数（原文 pp.19–21）
% ============================================================
\subsection{谱 Eisenstein 级数}

\pnumt{设定。}{}\ 我们有标准的对应：
\begin{equation*}
  \begin{tikzcd}[column sep=large, row sep=large]
    & \LS^{\restr}_{\bB} \arrow[dl, "q"'] \arrow[dr, "p"] & \\
    \LS^{\restr}_{\bT} & & \LS^{\restr}_{\bG}.
  \end{tikzcd}
\end{equation*}
映射 $p$ 可表且紧合，而映射 $q$ 拟光滑且 1-可表。此外，这些空间各自带有一个 Frobenius 自映射，且映射 $p$ 与 $q$ 缠绕这些 Frobenius。因此，由 (3.3) 与 (3.4)，我们得到典范映射：
\[
  \Gamma(\LS^{\arthm}_{\bT}, \omega_{\LS^{\arthm}_{\bT}}) \longrightarrow
  \Gamma(\LS^{\arthm}_{\bB}, \omega_{\LS^{\arthm}_{\bB}}) \longrightarrow
  \Gamma(\LS^{\arthm}_{\bG}, \omega_{\LS^{\arthm}_{\bG}}).
\]

\begin{defn}
上述映射的复合称为\emph{谱 Eisenstein 级数}
$\Eisspec : \Gamma(\LS^{\arthm}_{\bT}, \omega_{\LS^{\arthm}_{\bT}}) \to
\Gamma(\LS^{\arthm}_{\bG}, \omega_{\LS^{\arthm}_{\bG}})$。
\end{defn}

\subsection{谱 vs.\ 函数论 Eisenstein 级数}

\pnum{回忆：\emph{受限几何朗兰兹}（restricted geometric Langlands，见猜想 2.3.1）预测范畴之间存在如下等价：}
\begin{equation}
  \Shv_{\Nilp}(\Bun_G) \simeq \IndCoh_{\Nilp^{\mathrm{spec}}}(\LS^{\restr}_{\bG}).
\end{equation}

这一等价应当满足各种相容性。我们在此强调其中两个特别值得注意的：
\begin{itemize}
  \item（Hecke 相容性）等价 (3.7) 是 $\QCoh(\LS^{\restr}_{\bG})$-模范畴之间的等价；这里右端具有显然的作用，而左端带有 [2] 定理 14.3.2 的作用。
  \item（Eisenstein 相容性，$P = B$ 情形\zy{$P = B$ 指把抛物子群 $P$ 取成 Borel 子群 $B$ 的情形，即极小抛物情形。[来源：A/原文 §3.4.1 括注「$P = B$ case」]}）图表
  \begin{equation*}
    \begin{tikzcd}[column sep=large, row sep=large]
      \qLisse(\Bun_T) \arrow[r, "\Eis_!"] \arrow[d]
        & \Shv_{\Nilp}(\Bun_G) \arrow[d, "\cong"] \\
      \QCoh(\LS^{\restr}_{\bT}) \arrow[r, "\Eisspec"]
        & \IndCoh_{\Nilp^{\mathrm{spec}}}(\LS^{\restr}_{\bG})
    \end{tikzcd}
  \end{equation*}
  交换；这里左侧的竖直箭头是 [2] 例 21.2.9 中无条件构造出的等价。
\end{itemize}

\pnum{我们现在回忆如下结果：}

\begin{prop}
设 $\mathcal{Y}$ 是拟光滑的 ind-代数叠，带有自映射 $\phi_{\mathcal{Y}}$。令
$\mathcal{N} \subseteq T^*[-1]\mathcal{Y}$ 是闭锥形子叠（位于 $\mathcal{Y}$ 的 $-1$-平移余切丛中），使得对每个点
$y \in \mathcal{Y}^\phi$，映射 $d\phi[-1] : \mathcal{N}_y \to \mathcal{N}_y$ 都收缩到 $0 \in T^*_y[-1]\mathcal{Y}$。则映射
\[
  \tr_{\IndCoh_{\mathcal{N}}(\mathcal{Y})}(\phi^!)
  \longrightarrow \tr_{\IndCoh(\mathcal{Y})}(\phi^!)
  = \Gamma(\mathcal{Y}^\phi, \omega_{\mathcal{Y}^\phi})
\]
是同构。
\end{prop}

见 [2] §24.6.8。\footnote{在前引文献（loc.\ cit.）中，这一结果是以猜想的形式表述的。但它实际上可以由 [1] 的形式体系直接证明。}

特别地，我们发现受限几何朗兰兹给出同构：
\begin{equation}
  \Aut^{\unr}_{G,c} \simeq \Gamma(\LS^{\arthm}_{\bG}, \omega_{\LS^{\arthm}_{\bG}}).
\end{equation}

在受限几何朗兰兹的 Hecke 相容性之下，这是
$\Exc_{\bG} := \Gamma(\LS^{\arthm}_{\bG}, \OO_{\LS^{\arthm}_{\bG}})$-模之间的等价，其中左端从 [4] 继承其
$\Exc_{\bG}$-模结构。我们提醒（见 [2] §24.2）：$\Aut^{\unr}_c$ 上的 $\Exc_{\bG}$-模结构细化了
巡游算子在它上面的 Lafforgue--Xue 作用（见 §1.2.3）。

\pnum{我们现在得到：}

\begin{prop}
假设受限几何朗兰兹对 $G$ 成立，并带有其 Eisenstein 相容性。则如下图表交换：
\begin{equation*}
  \begin{tikzcd}[column sep=large, row sep=large]
    \Aut^{\unr}_{T,c} \arrow[r, "\text{ps-Eis}"] \arrow[d, "\cong"]
      & \Aut^{\unr}_{G,c} \arrow[d, "\cong"] \\
    \Gamma(\LS^{\arthm}_{\bT}, \omega_{\LS^{\arthm}_{\bT}}) \arrow[r, "\Eisspec"]
      & \Gamma(\LS^{\arthm}_{\bG}, \omega_{\LS^{\arthm}_{\bG}})
  \end{tikzcd}
\end{equation*}
其中竖直同构来自 (3.8)。
\end{prop}

事实上，这来自 ps-Eis 与 $\Eisspec$ 通过迹的实现，见 (3.2) 与 §3.2.6。
